% ============================================================
%  CHAPTER 7 -- Statistical Concepts Used in Machine Learning
%  Presenter 7
% ============================================================
\section{Statistical Concepts Used}

\begin{frame}{Maximum Likelihood Estimation (MLE)}
Given data $\mathcal{D} = \{x_1, \dots, x_N\}$ and parameters $\theta$:
\[
\hat{\theta}_{\text{MLE}} = \arg\max_{\theta}\; \prod_{i=1}^{N} P(x_i \mid \theta)
= \arg\max_{\theta}\; \sum_{i=1}^{N} \log P(x_i \mid \theta)
\]
\begin{itemize}
    \item Maximises the probability of observed data under the model.
    \item Equivalent to minimising the \textbf{negative log-likelihood} (NLL).
    \item Foundation of most supervised learning objectives.
    \item \textbf{Example:} MLE for Gaussian estimates $\hat{\mu} = \bar{x}$, $\hat{\sigma}^2 = \frac{1}{N}\sum(x_i - \bar{x})^2$.
\end{itemize}
\end{frame}

\begin{frame}{Maximum A Posteriori (MAP)}
Incorporates a prior $P(\theta)$:
\[
\hat{\theta}_{\text{MAP}} = \arg\max_{\theta}\; P(\mathcal{D} \mid \theta)\, P(\theta)
= \arg\max_{\theta}\; \left[\sum_{i=1}^{N} \log P(x_i \mid \theta) + \log P(\theta)\right]
\]
\begin{itemize}
    \item Gaussian prior on $\theta$ $\Rightarrow$ \textbf{L2 regularisation} (weight decay).
    \item Laplace prior on $\theta$ $\Rightarrow$ \textbf{L1 regularisation} (sparsity).
    \item MAP = MLE when prior is uniform.
    \item Bridges frequentist (MLE) and Bayesian approaches.
\end{itemize}
\end{frame}

\begin{frame}{Bias vs.\ Variance}
\begin{columns}[T]
\column{0.48\textwidth}
\begin{block}{Bias}
Error from overly simplistic model assumptions.\\
High bias $\Rightarrow$ \textbf{underfitting}.
\end{block}
\column{0.48\textwidth}
\begin{block}{Variance}
Error from sensitivity to training data fluctuations.\\
High variance $\Rightarrow$ \textbf{overfitting}.
\end{block}
\end{columns}
\vspace{6pt}
\[
\text{Expected Error} = \text{Bias}^2 + \text{Variance} + \text{Irreducible Noise}
\]
\begin{itemize}
    \item \textbf{Goal:} Find optimal model complexity that minimises total error.
    \item Cross-validation helps estimate this trade-off empirically.
\end{itemize}
\end{frame}

\begin{frame}{Sampling, Estimation \& Hypothesis Testing}
\begin{itemize}
    \item \textbf{Sampling:} Drawing representative subsets from a population.
    \begin{itemize}
        \item Mini-batch SGD samples from training data.
        \item Bootstrap sampling for confidence intervals.
    \end{itemize}
    \item \textbf{Estimation:} Using samples to infer population parameters.
    \begin{itemize}
        \item Point estimates (MLE, MAP) vs.\ interval estimates (confidence intervals).
    \end{itemize}
    \item \textbf{Hypothesis Testing:}
    \begin{itemize}
        \item $H_0$: null hypothesis; $H_1$: alternative hypothesis.
        \item $p$-value: probability of observing data as extreme under $H_0$.
        \item Used to compare model performances (e.g.\ paired $t$-test).
    \end{itemize}
\end{itemize}
\end{frame}

\begin{frame}{Overfitting and Underfitting}
\begin{itemize}
    \item \textbf{Overfitting:} Model memorises training data, fails on unseen data (high variance).
    \begin{itemize}
        \item Remedies: regularisation, dropout, early stopping, more data.
    \end{itemize}
    \item \textbf{Underfitting:} Model is too simple to capture patterns (high bias).
    \begin{itemize}
        \item Remedies: increase model capacity, add features, train longer.
    \end{itemize}
    \item \textbf{Detection:}
    \begin{itemize}
        \item Train loss $\ll$ val loss $\Rightarrow$ overfitting.
        \item Both losses high $\Rightarrow$ underfitting.
    \end{itemize}
\end{itemize}
\end{frame}

\begin{frame}{Applications in ML Algorithms}
\begin{columns}[T]
\column{0.48\textwidth}
\textbf{Logistic Regression}
\begin{itemize}
    \item Models $P(y=1 \mid x)$ via sigmoid.
    \item Trained by MLE (cross-entropy loss).
    \item L2 penalty = MAP with Gaussian prior.
\end{itemize}
\column{0.48\textwidth}
\textbf{Neural Networks}
\begin{itemize}
    \item Loss functions grounded in MLE.
    \item Dropout $\approx$ Bayesian approximation.
    \item Batch Norm uses batch statistics.
\end{itemize}
\end{columns}
\vspace{8pt}
\textbf{Model Evaluation:}
\begin{itemize}
    \item Cross-validation: unbiased estimate of generalisation error.
    \item Statistical tests to compare classifiers (McNemar's test, paired $t$-test).
    \item Confidence intervals on metrics (Precision, Recall, F1).
\end{itemize}
\end{frame}

\begin{frame}{Numerical Example: MLE for Gaussian}
\begin{exampleblock}{Problem}
Given data $\mathcal{D}=\{3, 5, 7, 5, 10\}$. Estimate $\hat{\mu}_{\text{MLE}}$ and $\hat{\sigma}^2_{\text{MLE}}$ assuming data follows a Gaussian distribution.
\end{exampleblock}
\textbf{Solution:}
\begin{align*}
\hat{\mu}_{\text{MLE}} &= \bar{x} = \frac{3+5+7+5+10}{5} = \frac{30}{5} = \mathbf{6.0}\\[6pt]
\hat{\sigma}^2_{\text{MLE}} &= \frac{1}{N}\sum_{i=1}^{N}(x_i - \hat{\mu})^2\\
&= \frac{(3-6)^2+(5-6)^2+(7-6)^2+(5-6)^2+(10-6)^2}{5}\\
&= \frac{9+1+1+1+16}{5} = \frac{28}{5} = \mathbf{5.6}
\end{align*}
$\hat{\sigma}_{\text{MLE}} = \sqrt{5.6} \approx \mathbf{2.366}$
\end{frame}

\begin{frame}{Numerical Example: Bias-Variance Decomposition}
\begin{exampleblock}{Problem}
True value: $y=10$. Three models are trained on different samples.\\
Model A predictions: $\{9.8,\;10.1,\;10.0\}$ \quad Model B predictions: $\{8.0,\;12.0,\;10.0\}$
\end{exampleblock}
\textbf{Solution:}
\begin{columns}[T]
\column{0.48\textwidth}
\textbf{Model A:}
\begin{align*}
\bar{\hat{y}} &= 9.967\\
\text{Bias} &= 9.967-10 = -0.033\\
\text{Bias}^2 &= \mathbf{0.001}\\
\text{Var} &= \mathbf{0.016}
\end{align*}
$\Rightarrow$ Low bias, low variance
\column{0.48\textwidth}
\textbf{Model B:}
\begin{align*}
\bar{\hat{y}} &= 10.0\\
\text{Bias} &= 10.0-10 = 0\\
\text{Bias}^2 &= \mathbf{0.0}\\
\text{Var} &= \mathbf{2.667}
\end{align*}
$\Rightarrow$ Zero bias, high variance
\end{columns}
\vspace{4pt}
\textbf{Insight:} Model A is preferred -- lower total error despite tiny bias.
\end{frame}

% --- Summary Slide ---
\begin{frame}{Summary}
\begin{enumerate}
    \item Probability \& statistics are the \textbf{mathematical foundation} of ML.
    \item Core concepts: sample spaces, Bayes' theorem, distributions, expected values.
    \item Statistical measures guide data preprocessing and feature engineering.
    \item Bayes' theorem and Bayesian thinking enable principled learning from data.
    \item MLE and MAP provide the theoretical basis for training ML models.
    \item Understanding bias--variance trade-off is key to building models that generalise.
    \item Every ML algorithm -- from logistic regression to deep learning -- relies on statistics.
\end{enumerate}
\end{frame}
